% Root and B full literal audits passed; numerical records are controls only.
\subsection{Three cross states opposite a corner of order at most four}

\begin{lemma}[A literal unicyclic graph correspondence]
\label{C:unicyclic-graph-correspondence}
Let a finite monoid $M$ act on a connected finite graph $\Gamma$ by
maps that send every oriented edge at unit speed to one oriented edge,
or collapse it to a vertex. Suppose $\dim H_1(\Gamma;\mathbb Q)=1$.
Then the cofiber of $\Gamma_{hM}\to BM$ cannot have type
$\mathbb CP^2$.
\end{lemma}
\begin{proof}
If the second rational homology of the cofiber is nonzero,
Lemma~\ref{C:connected-graph-cone-degree} forces every actor to
induce degree $+1$ on $H_1(\Gamma;\mathbb Q)$.
The graph has a unique circuit $C$, with, say, $l$ edges.
The image of its circuit word is a closed walk of length at most
$l$. Degree $+1$ requires a traversal of that circuit exactly
once; there is no length available for a tree detour or a
collapsed circuit edge. Thus $C$ is strictly invariant, and
every actor restricts to an orientation-preserving cellular
automorphism of $C$.

The inclusion $C\to\Gamma$ is an actual equivariant underlying
homotopy equivalence. It remains an equivalence on action bars
and their cofibers. Identify the equally subdivided circuit with
the unit circle. Its actor image is a group of rotations by
multiples of $2\pi/l$, so $z\mapsto z^{l+1}$ commutes strictly
with every actor. The same scalar-degree lemma makes all positive
rational products on the cofiber zero. This contradicts the
projective-plane square. The unit-speed hypothesis is used in
the closed-walk length argument.
\end{proof}

\begin{lemma}[Three states over a two-element right-zero minimal ideal]
\label{C:right-zero-minimal-two-three-states}
Suppose that a finite monoid $N$ has minimal ideal
$K=\{p,q\}$ with $ab=b$ for $a,b\in K$.
For every $(N,M)$-biset $S$ with $|S|\leq3$, a collage of type
$\mathbb CP^2$ yields an ordinary right $M$-act cone on three
states of that actual homotopy type.
There is no bound on the orders of $N$ and $M$.
\end{lemma}
\begin{proof}
Every $n\in N$ fixes each element of $K$ by left multiplication:
for $k,h\in K$, $nk=n(hk)=(nh)k=k$.
The left actions $P,Q$ of $p,q$ are retractions onto a common
subset $U$, and every $N$-actor fixes $U$ pointwise.
The right $M$-action preserves $U$ and commutes with both
retractions. The monoid $N$ has a right zero, so $BN$ is
contractible.

If $|U|=1$, its unique state is fixed by both monoids and the
common-fixed-state argument kills the degree-two products.
If $U=S$, the left action is trivial and the collage is actually
equivalent to the ordinary $S$-act cone.
If $P=Q$, the common map is an $N$-equivariant retraction:
$nP=P$ and $Pn=P$, the latter since $pn\in K$ and both elements
of $K$ have the same action. It is also $M$-equivariant.
The core collage is the ordinary $U$-act cone, and the relative
quotient is a suspension. Here $|U|\leq2$, whose ordinary act
has zero degree-two products. A one-dimensional degree-two
space on the original collage either pulls back from this
retract or lifts to its suspension quotient; its square is
zero in either case.

The sole remaining case has $U=\{u_0,u_1\}$ and
$S=U\sqcup\{z\}$, with $Pz=u_0$, $Qz=u_1$.
Let $Z=B(N\ltimes S)$ be the actual left action category.
It is connected and simply connected. To see the latter
directly, pass to its fundamental groupoid. At each $u_i$
all endomorphism labels become identities, since their
endomorphism monoid is $N$ and $BN$ is contractible.
For any arrow $n:z\to nz$, compose with the arrow labeled
$p$ from $nz$. The resulting label $pn$ lies in $K$.
If $nz=z$, distinctness of $Pz,Qz$ forces $pn=p$, so
$p_z n_z=p_z$ and $n_z$ is an identity in the groupoid.
If $nz=u_i$, the same equation identifies $n_z$ with
the unique distinguished arrow $p_z$ or $q_z$ to $u_i$.
Thus the groupoid has precisely the two tree arrows from
$z$ to $u_0,u_1$, their inverses, and their composites;
there is no nontrivial loop.

In particular $\widetilde H_0(Z;\mathbb Z)=H_1(Z;\mathbb Z)=0$.
The augmented cellular or simplicial chains of $Z$, tensored
with the free left $M$-bar resolution, compute the projection
$Z_{hM}\to BM$. Its cofiber has zero second homology:
the reduced coefficient complex starts in homology degree two,
and the cofiber suspends it once. Since $BN$ is contractible,
this is the original collage. It therefore cannot be a
projective plane. Only the actual ordinary-act case remains,
and one or two states in it are excluded by the finite-line
lemma.
\end{proof}

\begin{theorem}[All opposite corners of order at most four on three states]
\label{C:small-opposite-corner-three-states}
Let $|N|\leq4$ and $1\leq|S|\leq3$ in a triangular biset
category with corners $M,N$. If its classifying space has type
$\mathbb CP^2$, then an ordinary three-state right $M$-act cone
has that actual homotopy type. In particular $|M|\geq18$.
For each fixed $n\in\{1,2,3,4\}$, the least total arrow count
among such categories with $|N|=n$ is exactly $21+n$.
\end{theorem}
\begin{proof}
When the trivial right rational $N$-module is projective, apply
Lemma~\ref{C:three-state-exact-coinvariants}.
Otherwise the minimal Rees ideal has at least two columns.
The order bound forces its maximal group and row set to be
singletons: a two-row or nontrivial-group ideal with two
columns already has at least four elements, and its identity
would have to be outside it. Hence the minimal ideal is a
right-zero band of size two or three.
The size-two case is covered by the preceding lemma.
In the size-three case, $N$ is exactly that right-zero band
with identity adjoined.

For this last monoid, the common image $U$ of its three
minimal retractions is fixed pointwise by $N$.
The cases $|U|=1$, $U=S$, or all three retractions equal
are dealt with exactly as in the preceding proof.
Otherwise $S$ has two core states and one outside state.
The left action category has a natural graph model with
these three states as vertices and one edge from the outside
state to its image under each of the three right-zero labels.
Indeed the endomorphism categories at core states are copies
of the contractible $BN$; postcomposition fixes every
right-zero edge label. Collapsing these categories yields
the graph naturally for every commuting right action, by
the same nerve-coordinate construction as
Lemma~\ref{C:right-zero-action-graph}.
Both core states occur among the three edge endpoints, so
this is a connected three-vertex, three-edge graph.
Every right actor sends an edge at unit speed to the edge
with the same label, or collapses it to a core vertex.
Lemma~\ref{C:unicyclic-graph-correspondence} excludes this case.

The resulting ordinary act has three states, so the exact
three-state act-cone minimum gives $|M|\geq18$.
It follows that the original total is at least $18+3+n=21+n$.
For equality use the accepted twenty-two-arrow ordinary act
model, whose source has order eighteen and whose act has three
states. For $n\geq2$, replace its trivial sink by the
identity-adjoined right-zero band of size $n-1$, acting
trivially on the three states. Its classifying space is
contractible, so the natural edge-bar projection gives an
actual equivalence with the original act cone. For $n=1$
use that cone itself. These examples have exactly $21+n$
arrows.
\end{proof}

\begin{remark}[Exact small-corner inventory and a bounded control]
\label{C:small-opposite-corner-inventory}
The elementary Rees argument above is sufficient for the theorem.
As an independent bounded control,
\path{code/a_small_nonprojective_corners.py} expands the complete
own order-two through order-four monoid catalogs under opposition
and ordinary isomorphism. There are respectively $2,7,35$
ordinary classes, and $0,1,6$ nonprojective right-trivial cases.
The six order-four cases consist of the identity-adjoined
right-zero three-band and five extensions of a right-zero
two-band $K$ by an extra actor $x$ and the identity.
Their $(x^2,\tau_x)$ possibilities are
\[
 (1,\mathrm{id}),\quad(1,\mathrm{swap}),\quad
 (x,\mathrm{id}),\quad(x,\mathrm{constant}),\quad
 (k,\mathrm{constant}_k).
\]
The exact nonprojective tables and their Rees data are recorded in
\path{code/outputs/a_small_nonprojective_corners.json}.
A separate full three-state left-action control is
\path{code/a_four_corner_three_states.py}, with output
\path{code/outputs/a_four_corner_three_states.json}.
It exhausts all labeled actions and retains the counts and literal normalized
$d_1,d_2$ matrices in the two-core, distinct-retraction branch;
all twenty-four such instances for the two-band extensions
have $H_0=\mathbb Q$ and $H_1=0$. These calculations are
controls of the structural proof, not premises for it.
\end{remark}
